% HAND-WRITTEN fragment -- placed by Adc.tex.  NEW 2026-08-25.
%
% Source: ~/chips/castalia/simruns/adcchar_20260825/README.md S4, with
%   mc/mc_readiness.txt      the three readiness checks off the executed netlist
%   mc/mismatch_census.txt   0 devices at mismatchflag=1, 8 explicitly at 0,
%                            1128 with no mismatchflag parameter at all
%   mc/mc_variation_probes.txt   8 samples each: variations=mismatch -> 8 random
%                            vars, sigma(code) 0.000 and sigma(i_avdd) 0.000 uA;
%                            variations=process -> 14 vars, sigma(i_avdd)
%                            1.516 uA; variations=all -> 22 vars, reproduces
%                            `process' to every printed digit
%   mc/mc_results.txt        100 samples: code 509 at 1.25 V and 886 at 1.85 V
%                            in EVERY sample (sigma = 0); 131 or 132 at 0.65 V
%                            (mean 131.72, sigma 0.451).  I_AVDD 4.748 uA mean,
%                            sigma 0.542 (11.4 %); power 12.971 uW mean, sigma
%                            1.356 (10.5 %), range 10.487-16.357 uW.  READY at
%                            8.555 and 13.155 us in every sample.
%   mc/mc_rdiv/              15 samples, mismatchflag forced to 1 on the eight
%                            rppolywo of the comparator reference divider,
%                            variations=mismatch: sigma(code-zero) 0.821 mV =
%                            0.52 LSB, range 439.10-441.71 mV, code at 1.25 V
%                            508-510, sigma(slope) 0.078 codes/V (0.012 %).
% Run conditions: 100 samples, variations=all, section castalia_pm_25, nominal
% corner, 37 C, four spectre jobs of 25 sharing seed 1 (firstrun 1/26/51/76);
% 0 errors, 100/100 decoded.  Spectre reports 22 random variables, of which
% 5 of the 6 subckt-based mismatch parameters are non-referenced.
% Precedent CDAC matching benches in myshkin: Cap_DAC_2V5D_Mismatch_Test and
% Cap_DAC_MOM_Pcell_Mismatch_Test.  No cell or device names in the prose.

\subsubsection{Monte Carlo: What This Netlist Cannot Measure}\label{sss:adc-mc}

A 100-sample process-and-mismatch Monte Carlo of the extracted converter runs
clean and measures nothing about matching. Two of the three input levels return
the same code in all 100 samples; the third dithers between the two codes
adjacent to the boundary it already sits on. A mismatch-only control returns
\(\sigma = 0\) on every level and on the supply current, and the full run
reproduces the process-only run to every printed digit.

Three structures account for it, and none is fixable in post-processing.
Extraction flattens away the inline model wrappers in which the kit's
per-instance mismatch parameters are declared, so the converter's transistors
sit on raw model cards that the statistical block has no path to --- the
simulator reports five of its six mismatch parameters as non-referenced. The
capacitor array is extracted linear capacitors with no statistical model of any
kind, so the matching that actually sets a charge-redistribution converter's
linearity is not represented at all. The one remaining matched structure, the
resistor divider that sets the comparator's reference, ships with its mismatch
flag cleared, which multiplies the model's mismatch terms to zero.

\textbf{The \(\sigma\) on this converter's gain and offset are therefore
unmeasured, not small.} Nothing below should be read as a matching result. What
the run does establish is the analog bias spread, which is process and is real:
\SI{11.4}{\percent} on the analog supply current and \SI{10.5}{\percent} on
total power, \SIrange{10.5}{16.4}{\micro\watt}, with conversion timing identical
in every sample.

Forcing the reference divider's mismatch on bounds the one contribution the
netlist can express: \SI{0.821}{\milli\volt}, \(0.52\) LSB of \(1\sigma\) on the
converter zero, \(\pm 1\) code, and \SI{0.012}{\percent} on gain --- the divider
moves the comparison threshold, not the array's charge balance. That is small
against the three-code offset the nominal converter already carries and smaller
still than the channel's own loop offset. Measuring the array itself needs a
different bench: a schematic capacitor DAC built from the kit's statistical
metal-fringe capacitor device rather than from extracted linear elements. A
precedent bench for exactly that exists in the legacy library and has not been
run against this array.
